\section{Problem setting and notation}
\label{sec:setting}
% ============================================================================

We write $\Otilde(\cdot),\Omegatilde(\cdot),\Thetatilde(\cdot)$ to hide factors polynomial in
$\log d$. Table~\ref{tab:notation} fixes the notation used throughout.

\begin{table}[t]
\centering
\caption{Notation.}
\label{tab:notation}
\small
\begin{tabular}{ll}
\toprule
symbol & meaning \\
\midrule
$d$ & activation width (number of neurons) \\
$F$ & number of features, $F>d$ \\
$\Phi\in\R^{d\times F}$ & code; column $\phi_i$ is the direction of feature $i$ \\
$G\in\R^{F\times d}$ & linear readout \\
$M=G\Phi\in\R^{F\times F}$ & readout interface, $\operatorname{rank}(M)\le d$ \\
$D=\operatorname{diag}(M_{11},\dots,M_{FF})$ & diagonal gains; $D^{-1}M$ is the calibrated interface \\
$A=M-I_F$ & cross-talk matrix (zero diagonal after calibration) \\
$b\in\{0,1\}^F$ & sparse Boolean state; $S=\operatorname{supp}(b)$, $s=|S|$ \\
$\mu=\max_{i\ne j}|\langle\phi_i,\phi_j\rangle|$ & coherence of a unit-norm code \\
$W(F,d)=\frac{F-d}{d(F-1)}$ & mean-square Welch-type floor \\
$\widehat W$ & measured mean squared off-diagonal cross-talk \\
$\theta$ & decoding threshold (default $1/2$) \\
$p_{\mathrm{rec}}(s)$ & probability of exact threshold recovery at sparsity $s$ \\
$s_{95}$ & largest $s$ such that $p_{\mathrm{rec}}\ge0.95$ for \emph{every} tested $s'\le s$ \\
$\Sigma=\Phi\Phi^{\top}$ & frame operator (invertible iff $\Phi$ has full row rank) \\
$h_i=\phi_i^{\top}\Sigma^{-1}\phi_i$ & leverage of feature $i$; $\sum_ih_i=d$ \\
$W_{\star}(\Phi)$ & the code's own floor: least mean-square cross-talk it admits \\
$R_{\mathrm{geom}},R_{\mathrm{readout}}$ & $W_{\star}(\Phi)/W(F,d)$ and $W(G,\Phi)/W_{\star}(\Phi)$ \\
$\Phi_{-i}$ & $\Phi$ with column $i$ deleted \\
$\alpha\in(0,1]$ & smallest detectable active amplitude \\
$\kappa_i(\alpha)$ & collision frontier of feature $i$; separable at $s$ iff $\kappa_i>s$ \\
$\delta_i(s),\gamma_i(s)$ & distance between the two hulls, and the margin $\delta_i(s)/2$ \\
$m,m'$ & feature counts \emph{inside imported statements}; $m'$ plays the role of $F$ \\
\bottomrule
\end{tabular}
\end{table}

Two decoding questions are compared throughout. \emph{Analog reconstruction} asks that the
scores $z=Mb$ approximate $b$ in a norm; \emph{support recovery} asks only that
$\mathbf 1\{z_i\ge\theta\}=b_i$ for every $i$. The first is a statement about magnitudes and is
never exactly satisfiable when $F>d$; the second is a statement about a decision boundary and
can be satisfied exactly --- by a thresholded linear score as much as by anything else, which
is why that rule is measured here too rather than treating support recovery as the preserve of
nonlinearity. Everything below is an elaboration of that asymmetry.

% ============================================================================
